% ==================================================
%  Capítulo 2 — Marco Teórico
% ==================================================
\section{Marco Teórico}\label{sec:marco_teorico}

\subsection{Ecuaciones de Maxwell y Propagación Electromagnética}\label{subsec:maxwell}

El comportamiento de la luz al interactuar con nanoestructuras metálicas se describe mediante las ecuaciones de Maxwell, que gobiernan la evolución acoplada de los campos eléctrico $\mathbf{E}$ y magnético $\mathbf{H}$. En un medio lineal, isotrópico y no magnético, una forma diferencial útil para este trabajo es:

\begin{align}
    \nabla \times \mathbf{E} &= -\mu_0 \frac{\partial \mathbf{H}}{\partial t}, \label{eq:faraday} \\
    \nabla \times \mathbf{H} &= \varepsilon_0 \varepsilon_r \frac{\partial \mathbf{E}}{\partial t} + \mathbf{J}, \label{eq:ampere}
\end{align}

\noindent donde $\mu_0$ es la permeabilidad del vacío, $\varepsilon_0$ la permitividad del vacío, $\varepsilon_r$ la permitividad relativa del medio y $\mathbf{J}$ la densidad de corriente. Las ecuaciones complementarias de divergencia establecen que:

\begin{align}
    \nabla \cdot \mathbf{B} = 0,  \\
    \nabla \cdot \mathbf{D} = \rho,
\end{align}

\noindent con $\mathbf{B} = \mu_0 \mathbf{H}$ y $\mathbf{D} = \varepsilon_0 \varepsilon_r \mathbf{E}$ bajo la aproximación constitutiva anterior.

En materiales metálicos como el oro, la permitividad relativa $\varepsilon_r(\omega)$ depende de la frecuencia y es, en general, una magnitud compleja. Una primera aproximación para describir esta respuesta dispersiva es el modelo de Drude:

\begin{equation}
\varepsilon_r(\omega) = 1 - \frac{\omega_p^2}{\omega^2 + i\gamma\omega}, \label{eq:drude}
\end{equation}

\noindent donde $\omega_p$ es la frecuencia de plasma del metal y $\gamma$ es la tasa de amortiguamiento. En el rango visible, los metales nobles pueden presentar una parte real negativa de la permitividad, condición asociada a la existencia de modos plasmónicos superficiales, mientras que la parte imaginaria describe pérdidas ópticas por absorción. Esta respuesta electromagnética dispersiva es una de las bases físicas que permiten confinar y modular la luz mediante nanoestructuras metálicas.

\subsection{Plasmones Superficiales y Transmisión Óptica Extraordinaria}\label{subsec:plasmones}

Los plasmones superficiales polaritónicos (SPPs, por \textit{Surface Plasmon Polaritons}) son modos electromagnéticos acoplados a oscilaciones colectivas de los electrones libres en la interfaz entre un metal y un dieléctrico. Estos modos se propagan a lo largo de la interfaz y presentan confinamiento transversal del campo electromagnético. Para una interfaz plana metal--dieléctrico, su relación de dispersión puede escribirse como:

\begin{equation}
    k_{\text{SPP}} = \frac{\omega}{c} \sqrt{\frac{\varepsilon_m \, \varepsilon_d}{\varepsilon_m + \varepsilon_d}},
    \label{eq:spp_dispersion}
\end{equation}

\noindent donde $\varepsilon_m$ y $\varepsilon_d$ son las permitividades del metal y del dieléctrico, respectivamente, y $c$ es la velocidad de la luz en el vacío.

Cuando la luz incide sobre una película metálica perforada con nano-rendijas de dimensiones sub-longitud de onda, pueden aparecer fenómenos de transmisión resonante asociados al acoplamiento entre la luz incidente, los modos de la rendija y los plasmones superficiales de la película. Este comportamiento se conoce como \textit{Transmisión Óptica Extraordinaria (EOT)}~\cite{ebbesen1998eot}, y se caracteriza por una transmitancia que, bajo ciertas condiciones geométricas y espectrales, puede superar la esperada por modelos clásicos de difracción para aberturas sub-longitud de onda.


En el contexto de las metalentes plasmónicas, cada nano-rendija puede entenderse, de manera aproximada, como un elemento local que modifica la fase y la amplitud de la luz transmitida. El ancho de la rendija $w$, la separación centro a centro entre rendijas $c$, el espesor de la película y las propiedades ópticas del material determinan conjuntamente la fase $\phi$ y la transmitancia $T$ del campo emergente. Al variar estos parámetros geométricos a lo largo de la superficie, es posible aproximar un perfil de fase espacialmente variable que concentre o redirija la luz, emulando el comportamiento de una lente convergente en una estructura ultradelgada.

\subsection{Método de Diferencias Finitas en el Dominio del Tiempo (FDTD)}\label{subsec:fdtd}

El método FDTD es una técnica de simulación electromagnética que resuelve las ecuaciones de Maxwell discretizando tanto el espacio como el tiempo. Fue propuesto originalmente por Kane Yee en 1966~\cite{yee1966} y constituye uno de los enfoques más utilizados para modelar la interacción luz-materia en geometrías complejas. El método discretiza los campos en una malla espacial intercalada (celda de Yee) y avanza en pasos temporales que satisfacen la condición de estabilidad de Courant-Friedrichs-Lewy (CFL), lo que garantiza la convergencia numérica de la solución.

\subsubsection{Condiciones de Frontera Absorbentes}

Para simular un dominio espacial infinito con una malla finita, se emplean capas absorbentes en los bordes del dominio de simulación. En este proyecto se utilizan absorbedores (\texttt{mp.Absorber}) en lugar de las capas perfectamente acopladas (PML) clásicas, dado que los absorbedores ofrecen mayor estabilidad numérica para fuentes incidentes tipo haz gaussiano y estructuras metálicas que se extienden hasta los bordes del dominio.

\subsubsection{Implementación: Meep}

Las simulaciones de este proyecto se realizan con \textbf{Meep}~\cite{oskooi2010meep} (\textit{MIT Electromagnetic Equation Propagation}), un software libre que implementa el método FDTD con soporte para materiales dispersivos, fuentes gaussianas bidimensionales y exportación de campos en formato HDF5. En este trabajo, la simulación opera en un dominio bidimensional (2D) y se monitorea el campo eléctrico $E_y$ en el plano de propagación.

La configuración se evalúa bajo polarización TM, consistente con la excitación de modos plasmónicos en estructuras de nano-rendijas. En este tipo de geometría, la componente del campo eléctrico perpendicular a las rendijas es relevante para inducir acumulación de carga en las paredes metálicas, acoplar la luz incidente a modos confinados de la rendija y favorecer fenómenos asociados a la transmisión óptica extraordinaria. Así, la elección de la polarización no es arbitraria, sino que responde a la física dominante del sistema modelado.

\subsection{Surrogate Modeling}\label{subsec:surrogate}

El \textit{Surrogate Modeling} (modelado subrogado) consiste en construir un modelo estadístico o de aprendizaje automático que aproxime la relación entrada--salida de un simulador computacionalmente costoso. En el contexto de este proyecto, el modelo subrogado busca aprender un mapeo aproximado de la forma:

\begin{equation}
    f_{\theta}: \mathbf{x} \mapsto \hat{\mathbf{y}}, \qquad \mathbf{x} \in \mathbb{R}^d, \quad \hat{\mathbf{y}} \in \mathbb{R}^m,
    \label{eq:surrogate}
\end{equation}

\noindent donde $\mathbf{x}$ representa los parámetros geométricos de la metalente, principalmente los anchos de las nano-rendijas y la separación centro a centro, $\hat{\mathbf{y}}$ representa la respuesta óptica predicha, $\theta$ los parámetros aprendidos del modelo, y $d$ y $m$ las dimensionalidades de entrada y salida, respectivamente. En el enfoque final de este trabajo, $\hat{\mathbf{y}}$ corresponde al campo óptico focal complejo, representado mediante perfiles de amplitud y fase re-muestreados a longitud fija. A partir de esta respuesta predicha pueden calcularse métricas derivadas, como el coeficiente $R^2$ asociado al perfil de intensidad focal.

La ventaja fundamental de un modelo subrogado es que, una vez entrenado, su tiempo de inferencia puede ser del orden de milisegundos o menor, en contraste con los minutos u horas que puede requerir una evaluación mediante FDTD. Esto permite evaluar un gran número de configuraciones candidatas durante la etapa de búsqueda y reservar las simulaciones electromagnéticas directas para un conjunto reducido de candidatos seleccionados.

El diseño inverso asistido por modelos de aprendizaje es un área activa en nanofotónica~\cite{molesky2018inverse}. Entre los algoritmos candidatos para esta tarea se encuentran los Procesos Gaussianos (GP), los métodos de ensamble basados en \textit{Gradient Boosting} (XGBoost~\cite{chen2016xgboost}, LightGBM) y los Perceptrones Multicapa (MLP). La elección del algoritmo está condicionada por la naturaleza de la salida $\hat{\mathbf{y}}$. Cuando el objetivo es un escalar o un vector de baja dimensión, por ejemplo una figura de mérito que resume la calidad de la respuesta, los métodos de ensamble y los Procesos Gaussianos resultan alternativas naturales. Cuando, en cambio, el objetivo es un objeto estructurado de alta dimensión, como un espectro de transmisión o un campo electromagnético muestreado espacialmente, las arquitecturas neuronales con múltiples salidas resultan más adecuadas, pues pueden aprender de forma conjunta las correlaciones entre las componentes del vector de salida.

\subsection{Diseño Inverso y Optimización}\label{subsec:diseno_inverso}

El problema directo en el diseño de metalentes consiste en calcular la respuesta óptica dada una geometría específica, es decir, evaluar el mapeo $\mathbf{x} \rightarrow \mathbf{y}$. El diseño inverso plantea el problema complementario: dada una respuesta óptica deseada $\mathbf{y}^*$, por ejemplo un perfil de intensidad focal asociado a enfoque, se busca identificar una geometría $\mathbf{x}$ que produzca una respuesta suficientemente cercana a dicho objetivo.

En un esquema asistido por modelos subrogados, esta búsqueda puede formularse como:

\begin{equation}
\mathbf{x}^* = \arg\min_{\mathbf{x} \in \mathcal{X}}\; \mathcal{L}\!\left(f_{\theta}(\mathbf{x}),\; \mathbf{y}^*\right),
\label{eq:inverse_opt}
\end{equation}

\noindent donde $\mathcal{L}$ es una función de pérdida o discrepancia entre la respuesta predicha por el surrogate y el objetivo deseado, y $\mathcal{X}$ es el espacio de diseño factible definido por restricciones geométricas y físicas. Esta formulación debe interpretarse como una aproximación computacional al problema físico real, ya que $f_{\theta}$ no es el simulador exacto sino una representación aprendida a partir de datos.

Cuando la función objetivo es una caja negra no diferenciable, como ocurre al evaluar una respuesta predicha y luego post-procesarla mediante métricas de calidad, resultan apropiadas las metaheurísticas sin gradiente. Entre ellas, los algoritmos genéticos (GA) mantienen una población de soluciones candidatas y la hacen evolucionar mediante operadores de selección, cruzamiento y mutación. La función de aptitud puede combinar varios términos medibles, por ejemplo calidad de enfoque, concentración de intensidad o penalizaciones geométricas, cuando un solo escalar no discrimina suficientemente entre soluciones próximas. Las restricciones duras se incorporan reparando, proyectando o descartando candidatos fuera del espacio factible.

Una dificultad propia de la optimización asistida por surrogate es que el optimizador puede explotar errores del modelo. Si el surrogate sobreestima la calidad en regiones poco muestreadas, un optimizador libre puede converger hacia geometrías con alta aptitud predicha pero bajo desempeño real al ser evaluadas con FDTD. Este fenómeno es especialmente relevante en problemas de alta dimensión, con datos escasos y regiones óptimas estrechas. Por ello, la optimización sobre modelos subrogados requiere mecanismos de control, como validación periódica con el simulador, restricciones de búsqueda, penalización de extrapolación o ciclos activos de adquisición de nuevos datos.

En este contexto, resulta útil el concepto de \textit{manifold} o variedad empírica de soluciones. Aunque el espacio de diseño $\mathcal{X}$ pueda ser de alta dimensión, las geometrías observadas y físicamente relevantes suelen concentrarse en regiones de menor dimensión efectiva, determinadas por las familias parametrizadas del diseño de experimentos y por las restricciones físicas del sistema. Restringir la búsqueda a vecindarios de geometrías previamente observadas o validadas reduce el riesgo de extrapolación del surrogate y favorece una exploración más confiable de regiones prometedoras. Esta restricción no garantiza optimalidad global, pero puede mejorar la utilidad práctica del pipeline cuando el modelo subrogado no es suficientemente preciso en todo el espacio.

Por último, el \textit{hill-climbing} (ascenso de colina) es una estrategia de búsqueda local en la que se perturba iterativamente la mejor solución conocida y se conservan las modificaciones que mejoran la función objetivo. Aunque es un método simple y local, resulta útil para refinar candidatos y, cuando las perturbaciones se evalúan con el simulador FDTD, permite además sondear la estabilidad de una geometría frente a variaciones nanométricas.
